\section{Discussion}
\label{sec:discussion}

\subsection{Model World against Real World}

Figure~\ref{fig:modelreal} sets the two worlds side by side. For each risk measure it shows how many
times smaller AS makes it, first in the Monte Carlo (open diamonds) and then on real NSE minutes
(dots). The reductions are of the same order in both worlds: three to fourteen times, depending on
the measure. The model's promise about \emph{risk} is kept on real data even though several of its
assumptions bend or break. The reason is that the risk result needs only one assumption that does
hold: that a quote closer to the mid is filled more often (A1). Given that, skewing quotes towards
the side that unloads inventory must pull inventory back towards zero, whatever the price does.

\begin{figure}[!htbp]
\centering
\includegraphics[width=\textwidth]{summary_model_vs_reality}
\caption{How many times AS reduces each risk measure relative to the symmetric dealer, in the
model's world (open diamonds) and on real NSE minutes (dots). Log scale; 1 means no reduction.}
\label{fig:modelreal}
\end{figure}

The model's promise about \emph{profit} is not kept, because profit needs A6, no adverse selection,
and that is the assumption real order flow violates. In the Monte Carlo a fill carries no information
about the next price move, so the dealer keeps all the spread it captures. On the exchange, fills
cluster when prices move, as in the information models of Glosten and Milgrom~\cite{glosten1985}, and
the dealer gives the spread back. Inventory models explain why spreads are \emph{shaped} the way they
are; information explains why they must be wide enough to pay for adverse selection.

\subsection{The Market's Spread against the Model's}

Figure~\ref{fig:spreads} compares the spread the model would quote with the spread the market charges.
At midday the AS spread is \RELIANCESpreadMid{} rupees for RELIANCE and \KSCLSpreadMid{} for KSCL,
roughly three times the Roll~\cite{roll1984} spread of \RELIANCERoll{} and \KSCLRoll{}. Even the edge
term alone, which ignores inventory, is wider than the market's spread.

\begin{figure}[!htbp]
\centering
\includegraphics[width=\textwidth]{summary_market_vs_model_spread}
\caption{The market's spread (Roll estimator, and the average one-minute high--low range as an
upper reference) against the Avellaneda--Stoikov optimal spread at midday and at the open.}
\label{fig:spreads}
\end{figure}

This is the gap one expects between a single dealer and a competitive market. The model describes a
dealer who faces the whole fill curve $Ae^{-k\delta}$ alone and sets a price on it, like Garman's
monopolist~\cite{garman1976}. NSE is a limit order book with many liquidity providers, and
competition among them pushes the quoted spread down to about what adverse selection and inventory
costs require. A dealer who quoted the AS spread there would sit behind better prices and be filled
mainly when the price runs through the whole book, which is exactly the adverse-selection pattern
of Section~\ref{sec:replay}.

\subsection{Limitations}

\begin{enumerate}
\item \textbf{Bars, not the order book.} Fills are inferred from one-minute highs and lows. The proxy
  ignores queue position, so it is optimistic about how often a quote is filled. It also measures
  adverse selection over the same minute as the fill, which may overstate it compared with a
  tick-level simulation. Tick or order-book data would remove both problems.
\item \textbf{Short sample.} Yahoo keeps about 30 days of one-minute history, so the test window has
  nine days per stock. The risk results are large enough to be clear at this size; the mean-PnL
  results are not, as the KSCL counterexample shows.
\item \textbf{Frictions.} Quotes are not rounded to the tick, and there are no fees, securities
  transaction tax or exchange charges. Every one of these would lower PnL further, which strengthens
  the conclusion that the model is not a stand-alone source of profit.
\item \textbf{Constant parameters.} The quotes use one $\sigma$ for the whole day, which
  Section~\ref{sec:assumptions} shows is wrong by a factor of about three at the open.
\end{enumerate}

\subsection{Extensions}

Three extensions follow directly from the findings. First, a time-of-day volatility $\sigma(t)$ would
fix the assumption that breaks, widening the open and tightening midday. Second, the exact solution
of Gu\'eant, Lehalle and Fern\'andez-Tapia~\cite{gueant2013} removes the asymptotic approximation and
bounds inventory explicitly. Third, and most important, an adverse-selection term, for example making
fill intensity depend on recent order flow, would address the one assumption whose failure costs the
dealer money.
